\section{Discussion and Limitations}\label{sec:limitations}

Within its prespecified information geometry and working return model, the paper
establishes a bounded positive implication.
The target-anchored Wasserstein barycentric reconstruction maps
distribution-valued information positions into a predetermined, row-stochastic
barycentric interaction field, and the quadratic closure gives an economic
scale to the relation between stand-alone and peer-adjusted exposure in
\Cref{thm:spatial-closure}.
In the return panel, this field and the persistent news co-mention field select
largely different peers but generate correlated peer-return series.
Estimating them jointly leaves total model-implied adjustment close to its
single-field level and reallocates it across the two channels, a result that
separate one-field fits conceal.

The first limit is identification.
The quadratic criterion is an as-if reduced-form representation of exposure
adjustment, not a claim that managers observe the barycentric interaction field
or literally solve the displayed optimization problem.
Constructing every field from information ending in 2022 and freezing it before
the 2023--2026 return panel prevents evaluation returns from mechanically
determining their own peer weights.
That predetermination does not make text exogenous: industries, technologies,
investor attention, and state-dependent reporting selection can jointly shape
the information positions and returns.
The likelihood gains and channel decomposition therefore measure conditional
field fit rather than causal peer effects; isolating peer transmission would
require an exclusion restriction or shock design.
The full distinction between the structural adjustment parameter and the
pseudo-true target of QMLE remains the identification boundary in
\Cref{sec:identification-estimation}.

The second limit concerns measurement and design.
The central field contribution depends on a fixed representation and ground
metric for quadratic Wasserstein transport, balanced article clouds, and the
target-specific alignment and simplex restrictions of the target-anchored
Wasserstein barycentric reconstruction.
Balancing the clouds makes firms comparable but removes news volume as a source
of information, while changing the representation, ground metric, or
reconstruction design can alter the barycentric interaction field and its
conditional fit.
The implementation uses a predetermined deterministic rule to select one
optimal pairwise assignment, but some assignments are tied; the paper does not
establish that alternative optimal selections would leave $W^\flat$ unchanged.
Cross-geometry likelihood comparisons are consequently evidence about competing
measurement designs, not geometry-free rankings.
Prespecified alternatives using weighted or unbalanced clouds, other ground
metrics, and held-out evaluation windows would show which features of the
field survive those measurement choices.

Measurement also bounds the dispersion and multi-field conclusions.
The common carrier, its noncollapse constant $L$, and the transmission radii
$\tau_i$ in \Cref{eq:imported-transfer} remain latent and are not estimated by
the return QMLE.
The attenuation bracket therefore limits the share of implied stand-alone
exposure dispersion that peer adjustment can erase but does not determine its
level, so \Cref{cor:spatial-dispersion-certificate} remains a conditional
restriction rather than a measured quantity.
The two fields' peer-return series correlate at
\ppnum[3]{overlap-induced-pearson}, and the bootstrap correlation between
$\hat\rho_B$ and $\hat\rho_N$ is
\ppnum[3]{post-pooled-2023-2026-two-field-channel-correlation}.
Both channel intervals exclude zero in this sample, but more nearly collinear
fields could leave their division weakly determined even when total adjustment
remains precise.
The split should therefore be read from the joint region in
Figure~\ref{fig:two-field-region}; extensions to additional fields should report
the corresponding joint uncertainty and induced-signal collinearity.

The final limit is external validity.
The evidence comes from a large-cap, survivor-conditioned intersection of
\ppvalue{n-tickers-priced} firms with complete return histories, rather than a
representative cross-section of listed firms.
A broader or changing universe would alter each target's feasible peer set and
could change the field's stationary weighting.
The timing is also specific: the barycentric interaction field uses 2018--2022
text, the return evaluation begins in 2023, and the 2026 window ends on
\ppvalue{post-evaluation-date-end}, making that final period unsuitable for
like-for-like comparison with complete calendar years.
Reconstructing fields in successive ex ante windows and extending the sample to
firms with shorter histories would test whether the conditional patterns persist
across firm entry, market segments, and information regimes.
